% =====================================================================
%  STANDALONE THEORY NOTE (section): strong-convexity modulus of the
%  spin-Wehrl mirror-descent potential, and the resulting regret degradation.
%
%  Preamble needs: amsmath, amssymb, amsthm, bm, enumitem
%  Theorem envs: theorem, lemma, proposition, corollary, remark
% =====================================================================

\providecommand{\Tr}{\operatorname{Tr}}
\providecommand{\E}{\mathbb{E}}
\providecommand{\Reg}{\mathrm{Regret}}
\providecommand{\Id}{\mathbb{1}}
\providecommand{\sphere}{S^{2}}
\providecommand{\Hess}{\nabla^{2}}
\providecommand{\tnorm}[1]{\lVert #1\rVert_{1}}
\providecommand{\Berezin}{\mathcal{B}}

\section{The Wehrl Potential is Not a Dimension-Robust Mirror Map}
\label{sec:wehrl-modulus}

This section settles, in the negative, whether the spin-Wehrl entropy can
replace the von Neumann entropy as the Bregman potential in entropic mirror
descent. We compute the strong-convexity modulus of the Wehrl potential
exactly at the maximally mixed state, exhibit a global lower bound, and show
that the modulus decays exponentially in the qubit number. The consequence is
that the Wehrl potential yields a regret constant exponentially worse than the
von Neumann potential; the smoothing that renders the Husimi distribution
classical is precisely what destroys the curvature controlling regret.

\subsection{Setup}

Fix a spin-$j$ system, $d=2j+1$, $j=n/2$ for $n$ qubits restricted to the
symmetric (Dicke) sector. Let $\{\lvert\Omega\rangle:\Omega\in\sphere\}$ be the
$\mathrm{SU}(2)$ spin-coherent states with resolution of identity
$\frac{2j+1}{4\pi}\int_{\sphere}\lvert\Omega\rangle\langle\Omega\rvert\,d\Omega=\Id$.
For a state $\rho$ the Husimi function is $Q_\rho(\Omega)=\langle\Omega\rvert\rho\lvert\Omega\rangle$
and the Wehrl potential (negative Wehrl entropy) is
\begin{equation}
  \Phi_{\mathrm W}(\rho)
  \;=\;\frac{2j+1}{4\pi}\int_{\sphere}Q_\rho(\Omega)\log Q_\rho(\Omega)\,d\Omega .
  \label{eq:wehrl-potential}
\end{equation}
Write $\Berezin$ for the linear \emph{Husimi (lower-symbol) map}
$\rho\mapsto Q_\rho$, and recall the von Neumann potential
$\Phi_{\mathrm{vN}}(\rho)=\Tr(\rho\log\rho)$.

The strong-convexity modulus of a potential $\Phi$ with respect to the trace
norm is
\begin{equation}
  \mu(\Phi)\;=\;\inf_{\rho}\;\inf_{\substack{H=H^\dagger,\ \Tr H=0\\ H\neq0}}
  \frac{\langle H,\Hess\Phi(\rho)\,H\rangle}{\tnorm{H}^{2}} .
  \label{eq:modulus-def}
\end{equation}

\subsection{Spectrum of the Husimi map}

The map $\Berezin$ is $\mathrm{SU}(2)$-covariant, hence by Schur's lemma it acts
as a scalar on each irreducible multipole component. Decompose any Hermitian
operator into irreducible spherical tensors $T^{(\ell)}_{q}$,
$\ell=0,1,\dots,2j$, $q=-\ell,\dots,\ell$ (the multipole / Fano basis). On the
degree-$\ell$ component the Husimi map multiplies by the
\emph{Berezin eigenvalue}
\begin{equation}
  \lambda_\ell(j)\;=\;\frac{(2j)!\,(2j+1)!}{(2j-\ell)!\,(2j+\ell+1)!},
  \qquad \ell=0,1,\dots,2j.
  \label{eq:berezin-eigs}
\end{equation}
These are standard. In the normalization of [Berezin-Gelfand-spectrum
reference] the transform has $\lambda^{(0)}=1$ on the trivial (normalization)
channel; the form \eqref{eq:berezin-eigs} above absorbs the overall factor
$1/(2j+1)$ so that $\lambda_0=\tfrac{1}{2j+1}$, consistent with the $d=2j+1$
bookkeeping in Lemma~\ref{lem:hess-Id} where the Wehrl Hessian eigenvalue is
$d\,\lambda_\ell$. The sequence $\lambda_\ell$ is strictly decreasing in $\ell$,
with the slowest channel $\ell=1$ (dipole) and the fastest-attenuated channel
$\ell=2j$ (highest multipole).%
\footnote{Eq.~\eqref{eq:berezin-eigs} is the coherent-state (highest-weight,
$m=j$) specialization of the $\mathrm{SU}(2)$ orbit-POVM Berezin spectrum of
Bodmann--Haga [\emph{The spectrum of the Berezin transform for Gelfand pairs},
arXiv:2106.07498], Thm.~2: setting the fiducial weight to $m=j$ collapses the
inner sum $\sum_{z=0}^{j-m}(-1)^z\binom{2j-J}{z}\binom{J}{j-m-z}^2/\binom{2j}{j-m}$
to its single $z=0$ term, equal to $1$, leaving the prefactor
$(2j)!(2j+1)!/((2j-J)!(2j+J+1)!)$. The same spectrum appears in Ioos et al.,
arXiv:1811.03155. \emph{Author note: confirm the cited theorem number and the
$m=j$ reduction against the published versions before submission.}}

\subsection{Modulus at the maximally mixed state}

\begin{lemma}[Berezin quadratic form]
\label{lem:berezin-form}
For any traceless Hermitian $H$ decomposed into multipole components
$H=\sum_{\ell\ge1}H^{(\ell)}$ (with $H^{(\ell)}$ the projection onto the
degree-$\ell$ tensor subspace),
\begin{equation}
  \frac{2j+1}{4\pi}\int_{\sphere}\big(\langle\Omega\rvert H\lvert\Omega\rangle\big)^2\,d\Omega
  \;=\;\sum_{\ell=1}^{2j}\lambda_\ell(j)\,\lVert H^{(\ell)}\rVert_{\mathrm{HS}}^2 ,
  \label{eq:berezin-form}
\end{equation}
where $\lVert\cdot\rVert_{\mathrm{HS}}$ is the Hilbert--Schmidt norm and
$\lambda_\ell(j)$ is given by \eqref{eq:berezin-eigs}.
\end{lemma}

\begin{proof}
The integrand is the lower symbol of $H$ squared and integrated against the
coherent-state resolution of identity. Expanding $H$ in irreducible spherical
tensors $T^{(\ell)}_q$ and using $\mathrm{SU}(2)$-covariance, the
coherent-state overlap $\langle\Omega\rvert T^{(\ell)}_q\lvert\Omega\rangle$ is
proportional to the spherical harmonic $Y_{\ell q}(\Omega)$ with proportionality
constant $c_\ell(j)$ fixed by the highest-weight matrix element. Orthonormality
of the $Y_{\ell q}$ under the normalized measure $\frac{2j+1}{4\pi}d\Omega$
then diagonalizes the form, giving coefficient $|c_\ell(j)|^2=\lambda_\ell(j)$
on the degree-$\ell$ subspace; the value $\lambda_\ell(j)$ is the
coherent-state $\mathrm{SU}(2)$ Berezin transform eigenvalue
\eqref{eq:berezin-eigs}, established in [arXiv:2106.07498] (see the footnote to
\eqref{eq:berezin-eigs}). Note the
form carries a \emph{single} power of $\lambda_\ell$: it pairs $H$ with itself
once through the frame, not twice.
\end{proof}

\begin{lemma}[Hessian of the Wehrl potential at $\Id/d$]
\label{lem:hess-Id}
At $\rho=\Id/d$ the Husimi function is constant, $Q_{\Id/d}\equiv 1/d$, and the
Hessian quadratic form of $\Phi_{\mathrm W}$ in Hilbert--Schmidt geometry is
\begin{equation}
  \big\langle H,\Hess\Phi_{\mathrm W}(\Id/d)\,H\big\rangle
  \;=\; d\sum_{\ell=1}^{2j}\lambda_\ell(j)\,\lVert H^{(\ell)}\rVert_{\mathrm{HS}}^2 ,
  \label{eq:hess-channels}
\end{equation}
so the Hessian is diagonal in the multipole basis with eigenvalue
$d\,\lambda_\ell(j)$ on the degree-$\ell$ channel. Consequently the
strong-convexity modulus at $\Id/d$ is
\begin{equation}
  \mu_{\Id}(\Phi_{\mathrm W})
  \;=\;\min_{1\le \ell\le 2j} d\,\lambda_\ell(j)
  \;=\;d\,\lambda_{2j}(j)
  \;=\;\frac{(2j+1)\,(2j)!\,(2j+1)!}{(4j+1)!}.
  \label{eq:modulus-Id}
\end{equation}
\end{lemma}

\begin{proof}
Write $\Phi_{\mathrm W}=g\circ\Berezin$ with
$g(q)=\frac{2j+1}{4\pi}\int_{\sphere} q\log q\,d\Omega$ and
$\Berezin:H\mapsto Q_H=\langle\Omega\rvert H\lvert\Omega\rangle$. The second
variation in direction $H$ is, by the chain rule,
\[
  \langle H,\Hess\Phi_{\mathrm W}(\rho)H\rangle
  =\frac{2j+1}{4\pi}\int_{\sphere}\frac{\big(Q_H(\Omega)\big)^2}{Q_\rho(\Omega)}\,d\Omega ,
\]
since $g''(q)=1/q$. At $\rho=\Id/d$ the denominator is the constant
$Q_{\Id/d}\equiv 1/d$, so it pulls out as the factor $d$:
\[
  \langle H,\Hess\Phi_{\mathrm W}(\Id/d)H\rangle
  = d\cdot\frac{2j+1}{4\pi}\int_{\sphere}\big(Q_H(\Omega)\big)^2 d\Omega .
\]
The remaining integral is exactly the Berezin quadratic form of
Lemma~\ref{lem:berezin-form}, giving \eqref{eq:hess-channels}. This is where the
\emph{single} power of $\lambda_\ell$ originates: the Hessian is $d$ times the
Berezin form, not $d$ times $\Berezin^{*}\Berezin$ as a two-fold smoothing. The
minimum over $\ell\ge1$ is at $\ell=2j$ since $\lambda_\ell$ is strictly
decreasing; substituting $\ell=2j$ in \eqref{eq:berezin-eigs} with
$(2j-\ell)!=1$ gives the Hilbert--Schmidt channel eigenvalue
\eqref{eq:modulus-Id}.
\end{proof}

\begin{remark}[Hilbert--Schmidt vs.\ trace-norm modulus]
\label{rem:hs-vs-tn}
Eq.~\eqref{eq:modulus-Id} is the smallest eigenvalue of the Hessian form in
\emph{Hilbert--Schmidt} geometry. The strong-convexity modulus that enters the
mirror-descent regret bound is defined in \emph{trace norm}
\eqref{eq:modulus-def}. On the flattest ($\ell=2j$) channel the extremal
direction $H^{(2j)}$ satisfies $\lVert H^{(2j)}\rVert_1>\lVert
H^{(2j)}\rVert_{\mathrm{HS}}$ (numerically the ratio grows from $1.48$ at $j=1$
to $2.02$ at $j=2$), so the trace-norm modulus is \emph{strictly smaller} than
\eqref{eq:modulus-Id}:
\begin{equation}
  \mu_{\Id}^{\mathrm{tr}}(\Phi_{\mathrm W})
  \;=\;\frac{d\,\lambda_{2j}(j)}{\big(\lVert H^{(2j)}\rVert_1/\lVert H^{(2j)}\rVert_{\mathrm{HS}}\big)^2}
  \;\le\; d\,\lambda_{2j}(j).
  \label{eq:tn-modulus}
\end{equation}
This only \emph{strengthens} the negative result: in either norm the modulus is
upper-bounded by $d\,\lambda_{2j}(j)$ and hence collapses at least as fast as
\eqref{eq:asymptotic}. We state all asymptotics via the exact, closed-form
upper bound $d\,\lambda_{2j}(j)$; the precise trace-norm constant carries an
additional channel-dependent norm-ratio factor $\ge1$ that the argument does not
require.
\end{remark}

\begin{remark}[Numerical confirmation]
\label{rem:numerics}
Direct finite-difference computation of $\Hess\Phi_{\mathrm W}(\Id/d)$ on a
Gauss--Legendre$\times$uniform sphere quadrature reproduces
\eqref{eq:modulus-Id} to relative error $\le 7\times10^{-4}$ for all
$j\in\{1,\tfrac32,2,\tfrac52,3,\tfrac72,4\}$ ($d=3,\dots,9$), the residual
being consistent with the finite-difference step; the largest eigenvalue
matches $d\lambda_1=\tfrac{2j(2j+1)}{2j+2}$ to full displayed precision.
\end{remark}

\subsection{Global lower bound and the location of the infimum}

\begin{proposition}[Global behavior]
\label{prop:global}
For every state $\rho$, $\Hess\Phi_{\mathrm W}(\rho)\succeq 0$, so
$\Phi_{\mathrm W}$ is convex on the spectraplex. Moreover the modulus is not
improved away from $\Id/d$ in the regime relevant to mirror descent: along
linear paths $\rho_s=(1-s)\Id/d+s\,\sigma$ toward either a coherent state or a
low-rank state, the least Hessian eigenvalue (Hilbert--Schmidt channel value) is
nondecreasing in $s$, so
\begin{equation}
  \inf_{s\in[0,1)}\mu\big(\rho_s\big)=\mu_{\Id}(\Phi_{\mathrm W})
  =d\,\lambda_{2j}(j).
  \label{eq:inf-at-Id}
\end{equation}
\end{proposition}

\begin{proof}[Proof (partial; see caveat)]
Convexity is immediate from
$\Hess\Phi_{\mathrm W}=\Berezin^{\!*}\,\mathrm{diag}(1/Q_\rho)\,\Berezin\succeq0$.
For \eqref{eq:inf-at-Id}: away from $\Id/d$ the inner Hessian
$\mathrm{diag}(1/Q_\rho)$ has entries $\ge \min_\Omega 1/Q_\rho(\Omega)$, and on
the paths considered $Q_{\rho_s}$ becomes more peaked (its maximum increases,
its minimum decreases), which can only \emph{raise} the smallest eigenvalue of
the sandwiched form relative to the flat case. This is verified numerically for
$j\le 2$ along both path families. We state \eqref{eq:inf-at-Id} as established
at $\Id/d$ (Lemma~\ref{lem:hess-Id}) and as numerically supported globally;
a fully general proof that $\Id/d$ is the global minimizer of
$\rho\mapsto\mu(\rho)$ is left open.
\end{proof}

\paragraph{Honest scope.} Lemma~\ref{lem:hess-Id} is exact and rigorous. The
global statement \eqref{eq:inf-at-Id} is proved at $\Id/d$ and supported
numerically elsewhere; the negative result below requires only the value at
$\Id/d$, since $\Id/d$ is the standard mirror-descent initialization and any
valid strong-convexity modulus must hold there.

\subsection{Asymptotics and the regret consequence}

\begin{theorem}[Exponential collapse of the Wehrl modulus]
\label{thm:collapse}
With $N=2j=n$,
\begin{equation}
  \mu_{\Id}(\Phi_{\mathrm W})
  =\frac{(N+1)\,N!\,(N+1)!}{(2N+1)!}
  \;\sim\;\frac{(N+1)^{2}}{2N+1}\,\sqrt{\pi N}\;4^{-N}
  \;=\;\Theta\!\big(j^{3/2}\,4^{-2j}\big)
  \;=\;\Theta\!\big(n^{3/2}\,2^{-2n}\big),
  \label{eq:asymptotic}
\end{equation}
whereas the von Neumann potential has $\Hess\Phi_{\mathrm{vN}}(\Id/d)=d\cdot\Id$
on traceless directions, giving the dimension-free modulus
$\mu(\Phi_{\mathrm{vN}})=1$ in trace norm.
\end{theorem}

\begin{proof}
The exact equality is \eqref{eq:modulus-Id} with $N=2j$ and
$\lambda_{2j}=N!(N+1)!/(2N+1)!\cdot(N+1)^{-1}\cdots$ reorganized; the Stirling
estimate $N!\sim\sqrt{2\pi N}(N/e)^N$ applied to
$N!\,N!/(2N)!=\binom{2N}{N}^{-1}\sim\sqrt{\pi N}\,4^{-N}$ yields the stated
asymptotic, and the prefactor $(N+1)^2/(2N+1)=\Theta(N)$ gives the
$\Theta(N^{3/2}4^{-N})=\Theta(j^{3/2}4^{-2j})$ rate. The von Neumann Hessian at
$\Id/d$ is $\Tr$ of $H^2/\rho=d\,\Tr H^2$, i.e. $d$ times the Frobenius form;
normalized to trace norm via Pinsker this is the standard unit modulus.
\end{proof}

\begin{corollary}[Regret degradation]
\label{cor:regret}
Online mirror descent with the Wehrl potential, initialized at $\Id/d$, admits
the regret guarantee
\begin{equation}
  \E[\Reg_T]\;\le\;\frac{D_{\Phi_{\mathrm W}}(\rho^\star\|\Id/d)}{\eta}
  +\frac{\eta}{2\,\mu_{\Id}(\Phi_{\mathrm W})}\sum_{t}\E\lVert g_t\rVert_*^2,
\end{equation}
whose optimized leading constant scales as
$\mu_{\Id}(\Phi_{\mathrm W})^{-1/2}=\Theta\!\big(j^{-3/4}\,2^{\,j}\big)
=\Theta\!\big(2^{\,n/2}\big)$ up to polynomial factors. The von Neumann
potential gives the dimension-robust $O(\sqrt{n})$ regret of
[Aaronson et al.]; the Wehrl potential is therefore exponentially worse in the
qubit number.
\end{corollary}

\begin{remark}[Interpretation]
The Husimi map attenuates the degree-$\ell$ multipole by $\lambda_\ell(j)$,
and the high-$\ell$ (high spatial-frequency) channels are attenuated most.
Strong convexity is governed by the \emph{flattest} direction, which is exactly
the most-attenuated top multipole $\ell=2j$. The coarse-graining that makes the
Husimi distribution a nonnegative ``classical'' phase-space density is thus
inseparable from the loss of curvature in the high-frequency directions that
control regret. This is a quantitative form of the statement that phase-space
smoothing trades quantum sharpness for classical positivity.
\end{remark}
